\chapter[CONCLUSION]{Conclusion}
\label{ch:conclusion}

This proposal sets out a contactless framework for detecting corrosion on metal
surfaces by fusing RGB, thermal and LiDAR data, and for measuring the corroded
area and metal loss of each detected region from the LiDAR.

The case for the work rests on the limits of a single camera. Whether a corroded
component remains fit for service depends on how deep the metal loss goes as
well as how much area it covers, yet camera-based detection records appearance
only: it confuses stains, coatings and shadows with corrosion, cannot see
beneath intact paint, and reports severity as affected area. Thermal imaging
responds to defects below the surface and LiDAR measures surface geometry, but
the review in Chapter~\ref{ch:literature-review} found that, of 33
corrosion-sensing studies, 16 rely on RGB images alone, only one combines a
camera with thermal imaging, and none fuses RGB, thermal and LiDAR data. Nor has
any study measured what each of these sensors contributes to corrosion
detection under controlled conditions.

The research addresses this gap through three objectives: a pipeline that
registers the three sensors into a common reference frame; an early-fusion YOLO11
detector that localizes corroded regions in the registered five-channel data,
with the corroded area and metal loss of each region measured from the LiDAR;
and an ablation that compares each sensor alone, RGB with thermal, RGB with LiDAR
and all three, together with a blank-channel control, under identical conditions
and a decision rule fixed in advance.

A preliminary test on the public MID-1K dataset showed that the pipeline and the
ablation protocol run end to end. Early fusion reached a mean test mAP@0.5 of
0.596 against 0.572 for RGB alone, a gain too small to satisfy the decision
rule, and it matched the blank-channel control, so the depth channel added
nothing in that setup. The same thermal images fell from 0.756 to 0.521 mAP@0.5
when scored against RGB-aligned labels, which confirms that pixel-level
registration must precede fusion.

The expected contributions are a registered RGB--thermal--LiDAR corrosion
dataset, to be released if the owners of the scanned structures permit, and a
controlled measure of what each sensor contributes to corrosion detection,
including the case in which a sensor is not worth its cost. The main risk is
data acquisition; if field access is not secured within the schedule, the work
will be evaluated on controlled corrosion specimens and the resulting limit on
generalization will be reported.
